\section{Implementation notes}
\label{sec:implementation}

This section collects the points where an implementation has to go beyond the text of the
specification, or where the NIST reference implementation (NBIS) differs from it. Each item says
whether it was verified while writing this document. The section will grow with the Rust
implementation of the \texttt{wsq-viewer} project, which will add its choices of numerical
precision, its deviations from NBIS, if any, and measured performance.

\begin{description}[leftmargin=0pt]
  \item[Precision.] The specification's \emph{reference} transform is in double precision, while
    NBIS computes in single precision. The double-precision program used for
    Section~\ref{sec:compliance} differs from NBIS by at most one grey level, in fewer than
    0.01\% of the pixels, and by at most one quantizer index (verified).

  \item[The full tree.] NBIS splits the image into 60 subbands and never splits the quadrant of
    subbands 60--63, which encoder number two discards. A decoder that builds all 64 subbands
    produces the same image whenever $Q_{60..63} = 0$, and also decodes the (hypothetical) files
    that use them (verified on the NIST data, where all four are always zero).

  \item[All-ones Huffman codes.] All 86 WSQ files of the NIST decoder test set contain a table
    that assigns an all-ones code, which the specification forbids; NBIS decodes them with a
    warning. A decoder should accept them, and an encoder should never produce them (verified).

  \item[Variance switch.] The specification uses the window when the sum of the four variances
    exceeds 20\,000; NBIS uses it unless the sum is below 20\,000. The two differ only when the sum
    is exactly 20\,000 (from reading the NBIS source).

  \item[Small images.] With the other side large, an image needs at least 49 columns and 81 rows
    (81 pixels on each side for a square) for every variance window to hold two or more
    coefficients. Below that, Equation~\eqref{eq:variance} is undefined for some subband. The
    specification requires compliance only for images of at least 400 pixels on each side and
    says nothing about smaller ones, so an encoder has to choose a behaviour (verified by
    computing the windows for every size).

  \item[Constant images.] A uniform image has $R = 0$ and cannot be normalized by
    Equation~\eqref{eq:normalize}; an encoder has to treat it as a special case.

  \item[Blocks beyond three.] The block headers do not say which subbands a block holds. Since the
    indices are simply concatenated, a decoder decodes each block until its data ends and checks
    at the end that the total matches $\sum_{k:\,Q_k>0} X_kY_k$.

  \item[Untrusted input.] A viewer opens files from anywhere. Image dimensions come from the frame
    header and set the size of every allocation, symbol expansions can produce long runs, and
    tables may be missing, redefined or inconsistent. Every count read from the file has to be
    checked against what remains to be decoded before anything is written.

  \item[Coefficient range.] The escape codes carry at most 16-bit magnitudes, so an index must fit
    in $\pm 65535$. The authors of an independent Rust implementation report that on images with
    almost no detail, such as smooth gradients, encoder number two can produce bins fine enough to
    exceed this range, and that NBIS then wraps the values around \cite{wsqcodec}. This has not
    been verified here yet.
\end{description}
